% ============================================================================
%  gsd-command-reference.tex
%  GSD Command Reference — Wave 2 of the GSD User Guides Suite
%  Compile: xelatex (3-pass)
%  License: CC BY-SA 4.0 — Tibsfox
% ============================================================================
\documentclass[10pt,letterpaper]{article}

\title{GSD Command Reference}
\author{Tibsfox}
\date{}

\makeatletter
\def\input@path{{../shared/}}
\makeatother

% ============================================================================
%  gsd-guides-preamble.tex
%  Shared LaTeX preamble for the GSD User Guides Suite
%  Compile with: xelatex (3-pass), NOT pdflatex
%  License: CC BY-SA 4.0 — see \cclicense macro
%
%  Load order note: this file loads xcolor + tcolorbox BEFORE sourcing
%  gsd-color-scheme.tex, and then loads hyperref AFTER color definitions
%  are live. Documents only need to % ============================================================================
%  gsd-guides-preamble.tex
%  Shared LaTeX preamble for the GSD User Guides Suite
%  Compile with: xelatex (3-pass), NOT pdflatex
%  License: CC BY-SA 4.0 — see \cclicense macro
%
%  Load order note: this file loads xcolor + tcolorbox BEFORE sourcing
%  gsd-color-scheme.tex, and then loads hyperref AFTER color definitions
%  are live. Documents only need to % ============================================================================
%  gsd-guides-preamble.tex
%  Shared LaTeX preamble for the GSD User Guides Suite
%  Compile with: xelatex (3-pass), NOT pdflatex
%  License: CC BY-SA 4.0 — see \cclicense macro
%
%  Load order note: this file loads xcolor + tcolorbox BEFORE sourcing
%  gsd-color-scheme.tex, and then loads hyperref AFTER color definitions
%  are live. Documents only need to \input{gsd-guides-preamble} in their
%  preamble (this file pulls in the color scheme transitively).
% ============================================================================

% ---------- Fonts ----------
\usepackage{fontspec}
\setmainfont{DejaVu Serif}[
  Ligatures=TeX,
  Scale=0.95,
  BoldFont={DejaVu Serif Bold},
  ItalicFont={DejaVu Serif Italic},
  BoldItalicFont={DejaVu Serif Bold Italic}
]
\newfontfamily\headingfont{DejaVu Sans}[Scale=0.95]
\newfontfamily\monofont{DejaVu Sans Mono}[Scale=0.88]
\setmonofont{DejaVu Sans Mono}[Scale=0.88]

% ---------- Page geometry ----------
\usepackage[
  letterpaper,
  top=1in,
  bottom=1in,
  left=1.25in,
  right=1.25in,
  headheight=14pt
]{geometry}

% ---------- Color + callout engine (must load BEFORE color-scheme input) ----
\usepackage{xcolor}
\usepackage[most]{tcolorbox}

% ---------- Typesetting packages ----------
\usepackage{tabularx}
\usepackage{longtable}
\usepackage{booktabs}
\usepackage{colortbl}
\usepackage{lastpage}
\usepackage{enumitem}
\usepackage{listings}
\usepackage{tikz}
\usepackage{fancyhdr}
\usepackage{titlesec}
\usepackage{parskip}
\usepackage{microtype}

% ---------- Pull in the color scheme (defines gsdnavy, gsdaccent, callouts) -
\input{gsd-color-scheme.tex}

% ---------- Hyperref (needs colors already defined) ----------
\usepackage[
  colorlinks=true,
  linkcolor=gsdnavy,
  urlcolor=gsdaccent,
  citecolor=gsdnavy,
  bookmarksnumbered=true
]{hyperref}

% ---------- Header / footer ----------
\pagestyle{fancy}
\fancyhf{}
\fancyhead[L]{\small\headingfont\@title}
\fancyhead[R]{\small\headingfont\leftmark}
\fancyfoot[L]{\small\headingfont GSD v1.35.0 — CC BY-SA 4.0}
\fancyfoot[R]{\small\headingfont Page \thepage\ of \pageref*{LastPage}}
\renewcommand{\headrulewidth}{0.4pt}
\renewcommand{\footrulewidth}{0.4pt}

\fancypagestyle{plain}{%
  \fancyhf{}
  \fancyfoot[L]{\small\headingfont GSD v1.35.0 — CC BY-SA 4.0}
  \fancyfoot[R]{\small\headingfont Page \thepage\ of \pageref*{LastPage}}
  \renewcommand{\headrulewidth}{0pt}
  \renewcommand{\footrulewidth}{0.4pt}
}

% ---------- Chapter / section styling ----------
\titleformat{\chapter}[display]
  {\headingfont\huge\bfseries\color{gsdnavy}}
  {\chaptertitlename\ \thechapter}{20pt}{\Huge}
\titleformat{\section}
  {\headingfont\Large\bfseries\color{gsdnavy}}
  {\thesection}{1em}{}
\titleformat{\subsection}
  {\headingfont\large\bfseries\color{gsdnavy}}
  {\thesubsection}{1em}{}

% ---------- Listings (code blocks) ----------
\lstdefinestyle{gsdcode}{
  basicstyle=\ttfamily\small,
  backgroundcolor=\color{gsdcodebg},
  frame=leftline,
  framerule=2pt,
  rulecolor=\color{gsdaccent},
  xleftmargin=1em,
  xrightmargin=0.5em,
  aboveskip=0.8em,
  belowskip=0.8em,
  breaklines=true,
  breakatwhitespace=true,
  showstringspaces=false,
  columns=flexible,
  keepspaces=true
}
\lstset{style=gsdcode}

% ---------- Semantic macros ----------
\newcommand{\cclicense}{%
  \textit{Released under CC BY-SA 4.0.\\
  Authored by Tibsfox. GSD itself is MIT-licensed — see gsd-build/get-shit-done.}%
}

\newcommand{\chapterheading}[1]{{\headingfont\LARGE\bfseries\color{gsdnavy}#1\par\medskip}}
\newcommand{\sectionheading}[1]{{\headingfont\Large\bfseries\color{gsdnavy}#1\par\smallskip}}
 in their
%  preamble (this file pulls in the color scheme transitively).
% ============================================================================

% ---------- Fonts ----------
\usepackage{fontspec}
\setmainfont{DejaVu Serif}[
  Ligatures=TeX,
  Scale=0.95,
  BoldFont={DejaVu Serif Bold},
  ItalicFont={DejaVu Serif Italic},
  BoldItalicFont={DejaVu Serif Bold Italic}
]
\newfontfamily\headingfont{DejaVu Sans}[Scale=0.95]
\newfontfamily\monofont{DejaVu Sans Mono}[Scale=0.88]
\setmonofont{DejaVu Sans Mono}[Scale=0.88]

% ---------- Page geometry ----------
\usepackage[
  letterpaper,
  top=1in,
  bottom=1in,
  left=1.25in,
  right=1.25in,
  headheight=14pt
]{geometry}

% ---------- Color + callout engine (must load BEFORE color-scheme input) ----
\usepackage{xcolor}
\usepackage[most]{tcolorbox}

% ---------- Typesetting packages ----------
\usepackage{tabularx}
\usepackage{longtable}
\usepackage{booktabs}
\usepackage{colortbl}
\usepackage{lastpage}
\usepackage{enumitem}
\usepackage{listings}
\usepackage{tikz}
\usepackage{fancyhdr}
\usepackage{titlesec}
\usepackage{parskip}
\usepackage{microtype}

% ---------- Pull in the color scheme (defines gsdnavy, gsdaccent, callouts) -
% ============================================================================
%  gsd-color-scheme.tex
%  Color palette + tcolorbox environment definitions for all GSD User Guides
% ============================================================================

\definecolor{gsdnavy}{HTML}{1B2A4A}
\definecolor{gsdaccent}{HTML}{CB3837}
\definecolor{gsdcodebg}{HTML}{F5F5F5}
\definecolor{gsdcodeborder}{HTML}{CCCCCC}

\definecolor{gsdtipbg}{HTML}{E8F4FD}
\definecolor{gsdtipbar}{HTML}{2B6CB0}

\definecolor{gsdwarnbg}{HTML}{FFF3CD}
\definecolor{gsdwarnbar}{HTML}{B78D00}

\definecolor{gsddangerbg}{HTML}{F8D7DA}
\definecolor{gsddangerbar}{HTML}{B02A37}

\definecolor{gsdsuccessbg}{HTML}{D4EDDA}
\definecolor{gsdsuccessbar}{HTML}{198754}

\definecolor{gsdnotebg}{HTML}{F0F0F5}
\definecolor{gsdnotebar}{HTML}{495057}

% ---------- Callout boxes ----------
\tcbset{
  gsdcalloutbase/.style={
    enhanced,
    breakable,
    left=10pt,
    right=8pt,
    top=6pt,
    bottom=6pt,
    boxrule=0pt,
    frame hidden,
    borderline west={4pt}{0pt}{#1},
    colback=#1,
    fonttitle=\bfseries\headingfont,
    fontupper=\small,
    arc=2pt
  }
}

\newtcolorbox{gsdtip}[1][Tip]{%
  gsdcalloutbase=gsdtipbg,
  borderline west={4pt}{0pt}{gsdtipbar},
  colback=gsdtipbg,
  title={#1}
}

\newtcolorbox{gsdwarning}[1][Warning]{%
  gsdcalloutbase=gsdwarnbg,
  borderline west={4pt}{0pt}{gsdwarnbar},
  colback=gsdwarnbg,
  title={#1}
}

\newtcolorbox{gsdimportant}[1][Important]{%
  gsdcalloutbase=gsddangerbg,
  borderline west={4pt}{0pt}{gsddangerbar},
  colback=gsddangerbg,
  title={#1}
}

\newtcolorbox{gsdnote}[1][Note]{%
  gsdcalloutbase=gsdnotebg,
  borderline west={4pt}{0pt}{gsdnotebar},
  colback=gsdnotebg,
  title={#1}
}

\newtcolorbox{gsdsuccess}[1][Success]{%
  gsdcalloutbase=gsdsuccessbg,
  borderline west={4pt}{0pt}{gsdsuccessbar},
  colback=gsdsuccessbg,
  title={#1}
}


% ---------- Hyperref (needs colors already defined) ----------
\usepackage[
  colorlinks=true,
  linkcolor=gsdnavy,
  urlcolor=gsdaccent,
  citecolor=gsdnavy,
  bookmarksnumbered=true
]{hyperref}

% ---------- Header / footer ----------
\pagestyle{fancy}
\fancyhf{}
\fancyhead[L]{\small\headingfont\@title}
\fancyhead[R]{\small\headingfont\leftmark}
\fancyfoot[L]{\small\headingfont GSD v1.35.0 — CC BY-SA 4.0}
\fancyfoot[R]{\small\headingfont Page \thepage\ of \pageref*{LastPage}}
\renewcommand{\headrulewidth}{0.4pt}
\renewcommand{\footrulewidth}{0.4pt}

\fancypagestyle{plain}{%
  \fancyhf{}
  \fancyfoot[L]{\small\headingfont GSD v1.35.0 — CC BY-SA 4.0}
  \fancyfoot[R]{\small\headingfont Page \thepage\ of \pageref*{LastPage}}
  \renewcommand{\headrulewidth}{0pt}
  \renewcommand{\footrulewidth}{0.4pt}
}

% ---------- Chapter / section styling ----------
\titleformat{\chapter}[display]
  {\headingfont\huge\bfseries\color{gsdnavy}}
  {\chaptertitlename\ \thechapter}{20pt}{\Huge}
\titleformat{\section}
  {\headingfont\Large\bfseries\color{gsdnavy}}
  {\thesection}{1em}{}
\titleformat{\subsection}
  {\headingfont\large\bfseries\color{gsdnavy}}
  {\thesubsection}{1em}{}

% ---------- Listings (code blocks) ----------
\lstdefinestyle{gsdcode}{
  basicstyle=\ttfamily\small,
  backgroundcolor=\color{gsdcodebg},
  frame=leftline,
  framerule=2pt,
  rulecolor=\color{gsdaccent},
  xleftmargin=1em,
  xrightmargin=0.5em,
  aboveskip=0.8em,
  belowskip=0.8em,
  breaklines=true,
  breakatwhitespace=true,
  showstringspaces=false,
  columns=flexible,
  keepspaces=true
}
\lstset{style=gsdcode}

% ---------- Semantic macros ----------
\newcommand{\cclicense}{%
  \textit{Released under CC BY-SA 4.0.\\
  Authored by Tibsfox. GSD itself is MIT-licensed — see gsd-build/get-shit-done.}%
}

\newcommand{\chapterheading}[1]{{\headingfont\LARGE\bfseries\color{gsdnavy}#1\par\medskip}}
\newcommand{\sectionheading}[1]{{\headingfont\Large\bfseries\color{gsdnavy}#1\par\smallskip}}
 in their
%  preamble (this file pulls in the color scheme transitively).
% ============================================================================

% ---------- Fonts ----------
\usepackage{fontspec}
\setmainfont{DejaVu Serif}[
  Ligatures=TeX,
  Scale=0.95,
  BoldFont={DejaVu Serif Bold},
  ItalicFont={DejaVu Serif Italic},
  BoldItalicFont={DejaVu Serif Bold Italic}
]
\newfontfamily\headingfont{DejaVu Sans}[Scale=0.95]
\newfontfamily\monofont{DejaVu Sans Mono}[Scale=0.88]
\setmonofont{DejaVu Sans Mono}[Scale=0.88]

% ---------- Page geometry ----------
\usepackage[
  letterpaper,
  top=1in,
  bottom=1in,
  left=1.25in,
  right=1.25in,
  headheight=14pt
]{geometry}

% ---------- Color + callout engine (must load BEFORE color-scheme input) ----
\usepackage{xcolor}
\usepackage[most]{tcolorbox}

% ---------- Typesetting packages ----------
\usepackage{tabularx}
\usepackage{longtable}
\usepackage{booktabs}
\usepackage{colortbl}
\usepackage{lastpage}
\usepackage{enumitem}
\usepackage{listings}
\usepackage{tikz}
\usepackage{fancyhdr}
\usepackage{titlesec}
\usepackage{parskip}
\usepackage{microtype}

% ---------- Pull in the color scheme (defines gsdnavy, gsdaccent, callouts) -
% ============================================================================
%  gsd-color-scheme.tex
%  Color palette + tcolorbox environment definitions for all GSD User Guides
% ============================================================================

\definecolor{gsdnavy}{HTML}{1B2A4A}
\definecolor{gsdaccent}{HTML}{CB3837}
\definecolor{gsdcodebg}{HTML}{F5F5F5}
\definecolor{gsdcodeborder}{HTML}{CCCCCC}

\definecolor{gsdtipbg}{HTML}{E8F4FD}
\definecolor{gsdtipbar}{HTML}{2B6CB0}

\definecolor{gsdwarnbg}{HTML}{FFF3CD}
\definecolor{gsdwarnbar}{HTML}{B78D00}

\definecolor{gsddangerbg}{HTML}{F8D7DA}
\definecolor{gsddangerbar}{HTML}{B02A37}

\definecolor{gsdsuccessbg}{HTML}{D4EDDA}
\definecolor{gsdsuccessbar}{HTML}{198754}

\definecolor{gsdnotebg}{HTML}{F0F0F5}
\definecolor{gsdnotebar}{HTML}{495057}

% ---------- Callout boxes ----------
\tcbset{
  gsdcalloutbase/.style={
    enhanced,
    breakable,
    left=10pt,
    right=8pt,
    top=6pt,
    bottom=6pt,
    boxrule=0pt,
    frame hidden,
    borderline west={4pt}{0pt}{#1},
    colback=#1,
    fonttitle=\bfseries\headingfont,
    fontupper=\small,
    arc=2pt
  }
}

\newtcolorbox{gsdtip}[1][Tip]{%
  gsdcalloutbase=gsdtipbg,
  borderline west={4pt}{0pt}{gsdtipbar},
  colback=gsdtipbg,
  title={#1}
}

\newtcolorbox{gsdwarning}[1][Warning]{%
  gsdcalloutbase=gsdwarnbg,
  borderline west={4pt}{0pt}{gsdwarnbar},
  colback=gsdwarnbg,
  title={#1}
}

\newtcolorbox{gsdimportant}[1][Important]{%
  gsdcalloutbase=gsddangerbg,
  borderline west={4pt}{0pt}{gsddangerbar},
  colback=gsddangerbg,
  title={#1}
}

\newtcolorbox{gsdnote}[1][Note]{%
  gsdcalloutbase=gsdnotebg,
  borderline west={4pt}{0pt}{gsdnotebar},
  colback=gsdnotebg,
  title={#1}
}

\newtcolorbox{gsdsuccess}[1][Success]{%
  gsdcalloutbase=gsdsuccessbg,
  borderline west={4pt}{0pt}{gsdsuccessbar},
  colback=gsdsuccessbg,
  title={#1}
}


% ---------- Hyperref (needs colors already defined) ----------
\usepackage[
  colorlinks=true,
  linkcolor=gsdnavy,
  urlcolor=gsdaccent,
  citecolor=gsdnavy,
  bookmarksnumbered=true
]{hyperref}

% ---------- Header / footer ----------
\pagestyle{fancy}
\fancyhf{}
\fancyhead[L]{\small\headingfont\@title}
\fancyhead[R]{\small\headingfont\leftmark}
\fancyfoot[L]{\small\headingfont GSD v1.35.0 — CC BY-SA 4.0}
\fancyfoot[R]{\small\headingfont Page \thepage\ of \pageref*{LastPage}}
\renewcommand{\headrulewidth}{0.4pt}
\renewcommand{\footrulewidth}{0.4pt}

\fancypagestyle{plain}{%
  \fancyhf{}
  \fancyfoot[L]{\small\headingfont GSD v1.35.0 — CC BY-SA 4.0}
  \fancyfoot[R]{\small\headingfont Page \thepage\ of \pageref*{LastPage}}
  \renewcommand{\headrulewidth}{0pt}
  \renewcommand{\footrulewidth}{0.4pt}
}

% ---------- Chapter / section styling ----------
\titleformat{\chapter}[display]
  {\headingfont\huge\bfseries\color{gsdnavy}}
  {\chaptertitlename\ \thechapter}{20pt}{\Huge}
\titleformat{\section}
  {\headingfont\Large\bfseries\color{gsdnavy}}
  {\thesection}{1em}{}
\titleformat{\subsection}
  {\headingfont\large\bfseries\color{gsdnavy}}
  {\thesubsection}{1em}{}

% ---------- Listings (code blocks) ----------
\lstdefinestyle{gsdcode}{
  basicstyle=\ttfamily\small,
  backgroundcolor=\color{gsdcodebg},
  frame=leftline,
  framerule=2pt,
  rulecolor=\color{gsdaccent},
  xleftmargin=1em,
  xrightmargin=0.5em,
  aboveskip=0.8em,
  belowskip=0.8em,
  breaklines=true,
  breakatwhitespace=true,
  showstringspaces=false,
  columns=flexible,
  keepspaces=true
}
\lstset{style=gsdcode}

% ---------- Semantic macros ----------
\newcommand{\cclicense}{%
  \textit{Released under CC BY-SA 4.0.\\
  Authored by Tibsfox. GSD itself is MIT-licensed — see gsd-build/get-shit-done.}%
}

\newcommand{\chapterheading}[1]{{\headingfont\LARGE\bfseries\color{gsdnavy}#1\par\medskip}}
\newcommand{\sectionheading}[1]{{\headingfont\Large\bfseries\color{gsdnavy}#1\par\smallskip}}

% ============================================================================
%  gsd-command-data.tex
%  Rendering macros for GSD command names, flags, files
%  Command inventory verified against github.com/gsd-build/get-shit-done
%  v1.35.0 (commit 66a5f93, released 2026-04-10)
% ============================================================================

% \gsdcmd{command-name}  — renders a GSD command name in accent color monospace
\newcommand{\gsdcmd}[1]{{\ttfamily\color{gsdaccent}/gsd-#1}}

% \gsdcolon{command-name} — renders with colon syntax (skill form)
\newcommand{\gsdcolon}[1]{{\ttfamily\color{gsdaccent}/gsd:#1}}

% \gsdflag{flag}  — renders a CLI flag in muted monospace
\newcommand{\gsdflag}[1]{{\ttfamily\footnotesize\color{gsdnotebar}-{}-#1}}

% \gsdfile{filename}  — renders a filename / path in italic monospace
\newcommand{\gsdfile}[1]{{\ttfamily\itshape #1}}

% \gsdkbd{key}  — renders a keyboard key
\newcommand{\gsdkbd}[1]{{\ttfamily\footnotesize\color{gsdnavy}[#1]}}

% ---------------------------------------------------------------------------
%  GSD v1.35.0 Command Inventory (73 commands, 16 categories)
%  Source: /tmp/gsd-upstream/docs/COMMANDS.md
%
%  Core Workflow:
%    /gsd-new-project        /gsd-new-workspace       /gsd-list-workspaces
%    /gsd-remove-workspace   /gsd-discuss-phase       /gsd-ui-phase
%    /gsd-plan-phase         /gsd-execute-phase       /gsd-verify-work
%    /gsd-next               /gsd-session-report      /gsd-ship
%    /gsd-ui-review          /gsd-audit-uat           /gsd-audit-milestone
%    /gsd-complete-milestone /gsd-milestone-summary   /gsd-new-milestone
%
%  Phase Management:
%    /gsd-add-phase          /gsd-insert-phase        /gsd-remove-phase
%    /gsd-list-phase-assumptions   /gsd-analyze-dependencies
%    /gsd-plan-milestone-gaps      /gsd-research-phase
%    /gsd-validate-phase
%
%  Navigation:
%    /gsd-progress  /gsd-resume-work  /gsd-pause-work  /gsd-manager  /gsd-help
%
%  Utility:
%    /gsd-explore    /gsd-undo       /gsd-import      /gsd-from-gsd2
%    /gsd-quick      /gsd-autonomous /gsd-do          /gsd-note
%    /gsd-debug      /gsd-add-todo   /gsd-check-todos /gsd-add-tests
%    /gsd-stats      /gsd-profile-user  /gsd-health   /gsd-cleanup
%
%  Diagnostics:
%    /gsd-forensics  /gsd-extract-learnings
%
%  Workstream Management:
%    /gsd-workstreams
%
%  Configuration:
%    /gsd-settings   /gsd-set-profile
%
%  Brownfield:
%    /gsd-map-codebase  /gsd-onboard  /gsd-migrate
%
%  AI Integration:
%    /gsd-ai-integration-phase  /gsd-eval-review
%
%  Update:
%    /gsd-update  /gsd-scan
%
%  Code Quality:
%    /gsd-code-review  /gsd-code-review-fix  /gsd-secure-phase
%
%  Fast & Inline:
%    /gsd-fast  /gsd-quick  /gsd-inline
%
%  Backlog & Thread:
%    /gsd-add-backlog  /gsd-review-backlog  /gsd-thread
%
%  State Management:
%    /gsd-pr-branch  /gsd-reapply-patches  /gsd-plant-seed  /gsd-trace
%
%  Community:
%    /gsd-join-discord
% ---------------------------------------------------------------------------


% Denser margins — this is a reference sheet, not a narrative guide
\usepackage{geometry}
\geometry{letterpaper, top=0.75in, bottom=0.85in, left=0.9in, right=0.9in, headheight=14pt}

% Remap chapter-level commands from the guides to article sections so
% chapters don't eat a full page each. The source uses \chapter, \section,
% \section*, \commandentry.
\let\chapter\section
\let\oldsection\section
% In article class, \chapter no longer exists — we just aliased above.
% Section titles styled via titleformat in the preamble may need rebinding:
\titleformat{\section}
  {\headingfont\Large\bfseries\color{gsdnavy}}
  {\thesection}{1em}{}
\titleformat{\subsection}
  {\headingfont\large\bfseries\color{gsdnavy}}
  {\thesubsection}{0.8em}{}

% Tighten paragraph spacing for a dense reference
\setlength{\parskip}{2pt plus 1pt}
\setlength{\parindent}{0pt}

% ---------- Command-entry macro ----------
% Usage:
%   \commandentry{name}{synopsis}{syntax}{produces}{example}
% Any of produces / example may be empty ({}).
\newcommand{\commandentry}[5]{%
  \par\smallskip
  \noindent{\ttfamily\bfseries\color{gsdaccent}/gsd:#1}\quad{\small #2}\par
  \noindent\textbf{\footnotesize Syntax:} {\ttfamily\footnotesize #3}\par
  \ifx\relax#4\relax\else
    \noindent\textbf{\footnotesize Produces:} {\footnotesize #4}\par
  \fi
  \ifx\relax#5\relax\else
    \noindent\textbf{\footnotesize Ex:} {\ttfamily\footnotesize #5}\par
  \fi
}

% Flag table helper — use inside a commandentry or directly after
\newcommand{\flagtable}[1]{%
  \noindent\begin{tabularx}{\linewidth}{@{}l X@{}}
  \toprule
  \textbf{\footnotesize Flag} & \textbf{\footnotesize Description} \\
  \midrule
  #1
  \bottomrule
  \end{tabularx}\par\smallskip
}

\begin{document}

% ---------- Title page ----------
\begin{titlepage}
\centering
\vspace*{2cm}
{\headingfont\Huge\bfseries\color{gsdnavy} GSD\par}
\vspace{0.4cm}
{\headingfont\LARGE\bfseries\color{gsdnavy} Command Reference\par}
\vspace{0.8cm}
{\headingfont\large Every command, every flag, every file\par}
\vspace{3cm}
{\large Tibsfox\par}
\vspace{0.4cm}
{\small For GSD v1.35.0 \par}
\vspace{0.4cm}
{\small Audience: all GSD users — pin this to the wall\par}
\vfill
{\footnotesize \cclicense\par}
\end{titlepage}

\tableofcontents
\clearpage

% ============================================================================
\chapter*{How to Read This Reference}
\addcontentsline{toc}{chapter}{How to Read This Reference}

Commands appear in their canonical v1.35.0 skill form (\texttt{/gsd:command}). The older dash form (\texttt{/gsd-command}) remains a valid alias in every supported runtime. Argument and flag syntax follows standard POSIX conventions: angle brackets \texttt{<name>} mark required arguments, square brackets \texttt{[name]} mark optional ones, and ellipses \texttt{...} indicate a repeating list. Flag tables list only flags documented in the upstream \gsdfile{docs/COMMANDS.md} source for v1.35.0.

\medskip
\noindent\textbf{Runtime invocation prefixes.} GSD v1.35.0 supports fourteen host runtimes: Claude Code, OpenCode, Gemini CLI, Kilo, Codex, GitHub Copilot, Cursor, Windsurf, Antigravity, Augment, Trae, Qwen Code, CodeBuddy, and Cline. Most of them use \texttt{/gsd-*} at a prompt. Codex uses \texttt{\$gsd-*}. Skill-style invocation (\texttt{/gsd:*}) is available on the runtimes whose skill surfaces expose it. Pick whichever form your runtime accepts; the behavior is identical.

\clearpage

% ============================================================================
\chapter{Core Workflow}

\commandentry{new-project}
{Bootstrap a new project with deep context gathering — runs research, produces all planning artifacts, and writes CLAUDE.md.}
{/gsd:new-project [-{}-auto @file.md]}
{PROJECT.md, REQUIREMENTS.md, ROADMAP.md, STATE.md, config.json, research/, CLAUDE.md}
{/gsd:new-project -{}-auto @prd.md}

\flagtable{%
\gsdflag{auto @file.md} & Auto-extract from a document, skip interactive questions \\
}

\commandentry{new-workspace}
{Create an isolated multi-repo workspace with its own \gsdfile{.planning/} directory and repo copies (worktrees or clones).}
{/gsd:new-workspace -{}-name <n> [-{}-repos r1,r2] [-{}-path /t] [-{}-strategy worktree|clone] [-{}-branch <b>] [-{}-auto]}
{WORKSPACE.md, .planning/, repo copies}
{/gsd:new-workspace -{}-name feature-b -{}-repos hr-ui,ZeymoAPI}

\flagtable{%
\gsdflag{name <name>} & Workspace name (required) \\
\gsdflag{repos r1,r2} & Comma-separated repo paths or names \\
\gsdflag{path /target} & Target directory (default: \gsdfile{\textasciitilde/gsd-workspaces/<name>}) \\
\gsdflag{strategy worktree|clone} & Copy strategy (default: \texttt{worktree}) \\
\gsdflag{branch <name>} & Branch to check out (default: \texttt{workspace/<name>}) \\
\gsdflag{auto} & Skip interactive questions \\
}

\commandentry{list-workspaces}
{List active GSD workspaces and their status. Scans \gsdfile{\textasciitilde/gsd-workspaces/} for WORKSPACE.md manifests.}
{/gsd:list-workspaces}
{}
{/gsd:list-workspaces}

\commandentry{remove-workspace}
{Remove a workspace and clean up git worktrees. Refuses removal if any repo has uncommitted changes; requires name confirmation.}
{/gsd:remove-workspace <name>}
{}
{/gsd:remove-workspace feature-b}

\commandentry{discuss-phase}
{Capture implementation decisions before planning through a Socratic question loop.}
{/gsd:discuss-phase [N] [-{}-auto] [-{}-batch] [-{}-analyze] [-{}-power]}
{\{phase\}-CONTEXT.md, \{phase\}-DISCUSSION-LOG.md}
{/gsd:discuss-phase 2 -{}-analyze}

\flagtable{%
\gsdflag{auto} & Auto-select recommended defaults for all questions \\
\gsdflag{batch} & Group questions for batch intake instead of one-by-one \\
\gsdflag{analyze} & Add trade-off analysis during discussion \\
\gsdflag{power} & File-based bulk answers from a prepared answers file \\
}

\commandentry{ui-phase}
{Generate a UI design contract for a frontend phase before planning or execution.}
{/gsd:ui-phase [N]}
{\{phase\}-UI-SPEC.md}
{/gsd:ui-phase 2}

\commandentry{plan-phase}
{Research, plan, and verify a phase. The workhorse of the main workflow.}
{/gsd:plan-phase [N] [-{}-auto] [-{}-research|-{}-skip-research] [-{}-gaps] [-{}-skip-verify] [-{}-prd <f>] [-{}-reviews] [-{}-validate] [-{}-bounce|-{}-skip-bounce]}
{\{phase\}-RESEARCH.md, \{phase\}-\{N\}-PLAN.md, \{phase\}-VALIDATION.md}
{/gsd:plan-phase 1 -{}-bounce}

\flagtable{%
\gsdflag{auto} & Skip interactive confirmations \\
\gsdflag{research} & Force re-research even if RESEARCH.md exists \\
\gsdflag{skip-research} & Skip domain research step entirely \\
\gsdflag{gaps} & Gap closure mode (reads VERIFICATION.md, skips research) \\
\gsdflag{skip-verify} & Skip plan-checker verification loop \\
\gsdflag{prd <file>} & Use a PRD file instead of discuss-phase for context \\
\gsdflag{reviews} & Replan incorporating cross-AI review feedback from REVIEWS.md \\
\gsdflag{validate} & Run state validation before planning begins \\
\gsdflag{bounce} & Run external plan bounce validation after planning \\
\gsdflag{skip-bounce} & Skip plan bounce even if enabled in config \\
}

\commandentry{execute-phase}
{Execute all plans in a phase with wave-based parallelization, or run a single wave.}
{/gsd:execute-phase <N> [-{}-wave N] [-{}-validate] [-{}-cross-ai|-{}-no-cross-ai]}
{\{phase\}-\{N\}-SUMMARY.md per plan, commits, \{phase\}-VERIFICATION.md}
{/gsd:execute-phase 1 -{}-wave 2}

\flagtable{%
\gsdflag{wave N} & Execute only the Nth wave in the phase \\
\gsdflag{validate} & Run state validation before execution begins \\
\gsdflag{cross-ai} & Delegate execution to an external AI CLI \\
\gsdflag{no-cross-ai} & Force local execution even if cross-AI is enabled \\
}

\commandentry{verify-work}
{User acceptance testing with auto-diagnosis — runs after a phase has been executed.}
{/gsd:verify-work [N]}
{\{phase\}-UAT.md, fix plans if issues found}
{/gsd:verify-work 1}

\commandentry{next}
{Automatically advance to the next logical workflow step. Reads project state and dispatches to the right command (discuss, plan, execute, verify, or complete-milestone).}
{/gsd:next}
{}
{/gsd:next}

\commandentry{session-report}
{Generate a session report with work summary, outcomes, and estimated resource usage.}
{/gsd:session-report}
{.planning/reports/SESSION\_REPORT.md}
{/gsd:session-report}

\commandentry{ship}
{Create a pull request from completed phase work with an auto-generated PR body.}
{/gsd:ship [N|version] [-{}-draft]}
{GitHub PR, STATE.md updated}
{/gsd:ship 4 -{}-draft}

\flagtable{%
\gsdflag{draft} & Create as a draft PR \\
}

\commandentry{ui-review}
{Retroactive 6-pillar visual audit of an implemented frontend. Runs stand-alone — no active GSD project required.}
{/gsd:ui-review [N]}
{\{phase\}-UI-REVIEW.md, screenshots under .planning/ui-reviews/}
{/gsd:ui-review 3}

\commandentry{audit-uat}
{Cross-phase audit of all outstanding UAT and verification items, producing a categorized report with a human test plan.}
{/gsd:audit-uat}
{Categorized audit report}
{/gsd:audit-uat}

\commandentry{audit-milestone}
{Verify a milestone met its definition of done. Runs gap analysis over all executed phases.}
{/gsd:audit-milestone}
{Audit report with gap analysis}
{/gsd:audit-milestone}

\commandentry{complete-milestone}
{Archive a milestone and tag the release.}
{/gsd:complete-milestone}
{MILESTONES.md entry, git tag}
{/gsd:complete-milestone}

\commandentry{milestone-summary}
{Generate a comprehensive project summary from milestone artifacts (for team onboarding and review). Includes architecture decisions, phase breakdown, requirements coverage, tech debt, and a getting-started guide.}
{/gsd:milestone-summary [version]}
{.planning/reports/MILESTONE\_SUMMARY-v\{version\}.md}
{/gsd:milestone-summary v1.0}

\commandentry{new-milestone}
{Start the next version cycle. Updates PROJECT.md and creates a new REQUIREMENTS.md and ROADMAP.md.}
{/gsd:new-milestone [name] [-{}-reset-phase-numbers]}
{Updated PROJECT.md, new REQUIREMENTS.md, new ROADMAP.md}
{/gsd:new-milestone "v2.0 Mobile"}

\flagtable{%
\gsdflag{reset-phase-numbers} & Restart the new milestone at Phase 1 and archive old phase dirs \\
}

% ============================================================================
\chapter{Phase Management}

\commandentry{add-phase}
{Append a new phase to the roadmap. Interactive — describe the phase to the prompter.}
{/gsd:add-phase}
{Updated ROADMAP.md}
{/gsd:add-phase}

\commandentry{insert-phase}
{Insert urgent work between existing phases using decimal numbering (3 becomes 3.1, etc.).}
{/gsd:insert-phase [N]}
{Updated ROADMAP.md with decimal phase}
{/gsd:insert-phase 3}

\commandentry{remove-phase}
{Remove a future phase and renumber subsequent phases.}
{/gsd:remove-phase [N]}
{Updated ROADMAP.md, renumbered phase dirs}
{/gsd:remove-phase 7}

\commandentry{list-phase-assumptions}
{Preview the assumed approach for a phase before committing to plan-phase work.}
{/gsd:list-phase-assumptions [N]}
{}
{/gsd:list-phase-assumptions 2}

\commandentry{analyze-dependencies}
{Analyze phase dependencies and suggest \texttt{Depends on} entries for ROADMAP.md. Run this before \gsdcolon{manager} when phases have empty \texttt{Depends on} fields.}
{/gsd:analyze-dependencies}
{Dependency suggestion table; optionally updates ROADMAP.md}
{/gsd:analyze-dependencies}

\commandentry{plan-milestone-gaps}
{Create phases to close gaps surfaced by a milestone audit.}
{/gsd:plan-milestone-gaps}
{New phases in ROADMAP.md}
{/gsd:plan-milestone-gaps}

\commandentry{research-phase}
{Deep ecosystem research only — standalone entry point. Usually run \gsdcolon{plan-phase} instead, which invokes research as its first step.}
{/gsd:research-phase [N]}
{\{phase\}-RESEARCH.md}
{/gsd:research-phase 4}

\commandentry{validate-phase}
{Retroactively audit and fill Nyquist validation gaps for a phase.}
{/gsd:validate-phase [N]}
{Updated \{phase\}-VALIDATION.md}
{/gsd:validate-phase 2}

% ============================================================================
\chapter{Navigation}

\commandentry{progress}
{Show project status and next recommended steps.}
{/gsd:progress}
{}
{/gsd:progress}

\commandentry{resume-work}
{Restore full session context after a context reset or when starting a new session.}
{/gsd:resume-work}
{}
{/gsd:resume-work}

\commandentry{pause-work}
{Save a context handoff when stopping mid-phase. Creates \gsdfile{continue-here.md}.}
{/gsd:pause-work}
{continue-here.md}
{/gsd:pause-work}

\commandentry{manager}
{Interactive command center for managing multiple phases from one terminal. Dispatches discuss inline; dispatches plan and execute as background agents.}
{/gsd:manager}
{}
{/gsd:manager}

\begin{gsdnote}
The manager respects the \texttt{manager.flags.discuss}, \texttt{manager.flags.plan}, and \texttt{manager.flags.execute} passthrough flags in \gsdfile{.planning/config.json} — see the Configuration Reference chapter.
\end{gsdnote}

\commandentry{help}
{Show all commands and a quick usage guide.}
{/gsd:help}
{}
{/gsd:help}

% ============================================================================
\chapter{Utilities}

\commandentry{explore}
{Socratic ideation session — guides an idea through probing questions, optionally spawns research, then routes output to the right GSD artifact (notes, todos, seeds, research questions, requirements, or a new phase).}
{/gsd:explore [topic]}
{Routed output (varies by destination)}
{/gsd:explore authentication strategy}

\commandentry{undo}
{Safe git revert for GSD phase or plan commits. Uses the phase manifest for dependency checks and always surfaces a confirmation gate before reverting.}
{/gsd:undo (-{}-last N | -{}-phase NN | -{}-plan NN-MM)}
{Revert commits}
{/gsd:undo -{}-phase 03}

\flagtable{%
\gsdflag{last N} & Show recent GSD commits for interactive selection \\
\gsdflag{phase NN} & Revert all commits for a phase \\
\gsdflag{plan NN-MM} & Revert all commits for a specific plan \\
}

\commandentry{import}
{Ingest an external plan file into the GSD planning system. Detects conflicts against PROJECT.md decisions before writing.}
{/gsd:import -{}-from <filepath>}
{New PLAN.md validated by gsd-plan-checker}
{/gsd:import -{}-from /tmp/team-plan.md}

\flagtable{%
\gsdflag{from <filepath>} & Path to the external plan file (required) \\
}

\commandentry{from-gsd2}
{Reverse migration from GSD-2 format (\gsdfile{.gsd/}) back to v1 \gsdfile{.planning/} format. Flattens Milestone/Slice hierarchy to sequential phase numbers.}
{/gsd:from-gsd2 [-{}-dry-run] [-{}-force] [-{}-path <dir>]}
{PROJECT.md, REQUIREMENTS.md, ROADMAP.md, STATE.md, sequential phase dirs}
{/gsd:from-gsd2 -{}-dry-run}

\flagtable{%
\gsdflag{dry-run} & Preview migration without writing \\
\gsdflag{force} & Overwrite existing .planning/ directory \\
\gsdflag{path <dir>} & Specify the GSD-2 root (default: current dir) \\
}

\commandentry{quick}
{Execute an ad-hoc task with GSD guarantees — composable flags turn quick-tasks into full phases on demand.}
{/gsd:quick [-{}-full] [-{}-discuss] [-{}-research]}
{.planning/quick/NNN-slug/ with PLAN.md and SUMMARY.md}
{/gsd:quick -{}-discuss -{}-research -{}-full}

\flagtable{%
\gsdflag{full} & Enable plan checking (2 iterations) and post-execution verification \\
\gsdflag{discuss} & Lightweight pre-planning discussion \\
\gsdflag{research} & Spawn a focused researcher before planning \\
}

\commandentry{autonomous}
{Run all remaining phases autonomously. Respects \texttt{workflow.skip\_discuss} to bypass the discuss step.}
{/gsd:autonomous [-{}-from N] [-{}-to N] [-{}-interactive]}
{Phase artifacts and commits for every executed phase}
{/gsd:autonomous -{}-from 3 -{}-to 5}

\flagtable{%
\gsdflag{from N} & Start from a specific phase number \\
\gsdflag{to N} & Stop after completing a specific phase number \\
\gsdflag{interactive} & Lean context with periodic user input checkpoints \\
}

\commandentry{do}
{Route freeform text to the right GSD command. After invocation, describe what you want.}
{/gsd:do}
{}
{/gsd:do}

\commandentry{note}
{Zero-friction idea capture. Append, list, or promote notes to todos.}
{/gsd:note [text | list | promote N] [-{}-global]}
{Project or global notes file}
{/gsd:note "Consider caching for API responses"}

\flagtable{%
\gsdflag{global} & Use global scope for the note operation \\
}

\commandentry{debug}
{Systematic debugging with persistent cross-session state. Subcommands manage active debug sessions.}
{/gsd:debug [-{}-diagnose] <description> \textbar{} list \textbar{} status <slug> \textbar{} continue <slug>}
{Debug session artifacts with Evidence, Eliminated, Resolution, TDD checkpoint}
{/gsd:debug "Login button not responding on mobile Safari"}

\flagtable{%
\gsdflag{diagnose} & Diagnosis-only mode — investigate without attempting fixes \\
}

\commandentry{add-todo}
{Capture an idea or task for later.}
{/gsd:add-todo [description]}
{todos/ entry}
{/gsd:add-todo "Consider adding dark mode support"}

\commandentry{check-todos}
{List pending todos and select one to work on.}
{/gsd:check-todos}
{}
{/gsd:check-todos}

\commandentry{add-tests}
{Generate tests for a completed phase.}
{/gsd:add-tests [N]}
{Test files committed atomically}
{/gsd:add-tests 2}

\commandentry{stats}
{Display project statistics dashboard.}
{/gsd:stats}
{}
{/gsd:stats}

\commandentry{profile-user}
{Generate a developer behavioral profile from Claude Code session analysis across 8 dimensions. Produces artifacts that personalize subsequent GSD responses.}
{/gsd:profile-user [-{}-questionnaire] [-{}-refresh]}
{USER-PROFILE.md, /gsd-dev-preferences command, CLAUDE.md profile section}
{/gsd:profile-user -{}-questionnaire}

\flagtable{%
\gsdflag{questionnaire} & Use interactive questionnaire instead of session analysis \\
\gsdflag{refresh} & Re-analyze sessions and regenerate the profile \\
}

\commandentry{health}
{Validate \gsdfile{.planning/} directory integrity. With \gsdflag{repair}, auto-fix recoverable issues.}
{/gsd:health [-{}-repair]}
{Integrity report, optional repair actions}
{/gsd:health -{}-repair}

\flagtable{%
\gsdflag{repair} & Auto-fix recoverable integrity issues \\
}

\commandentry{cleanup}
{Archive accumulated phase directories from completed milestones.}
{/gsd:cleanup}
{Archived phase dirs}
{/gsd:cleanup}

% ============================================================================
\chapter{Diagnostics}

\commandentry{forensics}
{Post-mortem investigation of failed or stuck GSD workflows. Checks git history, artifact integrity, STATE.md anomalies, and uncommitted work; covers at least 4 anomaly types.}
{/gsd:forensics [description]}
{.planning/forensics/report-\{timestamp\}.md}
{/gsd:forensics "Phase 3 execution stalled"}

\commandentry{extract-learnings}
{Extract reusable patterns, anti-patterns, and architectural decisions from completed phase work.}
{/gsd:extract-learnings <N> [-{}-all] [-{}-format markdown|json]}
{.planning/learnings/\{phase\}-LEARNINGS.md}
{/gsd:extract-learnings 3}

\flagtable{%
\gsdflag{all} & Extract learnings from all completed phases \\
\gsdflag{format} & Output format: \texttt{markdown} (default) or \texttt{json} \\
}

% ============================================================================
\chapter{Workstream Management}

\commandentry{workstreams}
{Manage parallel workstreams for concurrent work on different milestone areas. Subcommands select the action.}
{/gsd:workstreams [list | create <n> | status <n> | switch <n> | progress | complete <n> | resume <n>]}
{Workstream directories under .planning/, per-workstream state tracking}
{/gsd:workstreams create backend-api}

\begin{tabularx}{\linewidth}{@{}l X@{}}
\toprule
\textbf{\footnotesize Subcommand} & \textbf{\footnotesize Description} \\
\midrule
\texttt{list} & List all workstreams with status (default if no subcommand) \\
\texttt{create <name>} & Create a new workstream \\
\texttt{status <name>} & Detailed status for one workstream \\
\texttt{switch <name>} & Set active workstream \\
\texttt{progress} & Progress summary across all workstreams \\
\texttt{complete <name>} & Archive a completed workstream \\
\texttt{resume <name>} & Resume work in a workstream \\
\bottomrule
\end{tabularx}
\smallskip

% ============================================================================
\chapter{Configuration Commands}

\commandentry{settings}
{Interactive configuration of workflow toggles and model profile.}
{/gsd:settings}
{Updated .planning/config.json}
{/gsd:settings}

\commandentry{set-profile}
{Quick switch of the model profile.}
{/gsd:set-profile <profile>}
{Updated .planning/config.json}
{/gsd:set-profile budget}

\flagtable{%
\texttt{quality} & Opus for all decision-making, Sonnet for verification \\
\texttt{balanced} & Opus for planning, Sonnet for everything else (default) \\
\texttt{budget} & Sonnet for code-writing, Haiku for research and verification \\
\texttt{inherit} & All agents use the current session model \\
}

% ============================================================================
\chapter{Brownfield \& Analysis}

\commandentry{map-codebase}
{Analyze an existing codebase with four parallel mapper agents. Use before \gsdcolon{new-project} on a brownfield repo.}
{/gsd:map-codebase [area]}
{.planning/codebase/ analysis files}
{/gsd:map-codebase auth}

\commandentry{scan}
{Rapid single-focus codebase assessment — lightweight alternative to map-codebase (one agent instead of four).}
{/gsd:scan [-{}-focus tech|arch|quality|concerns|tech+arch]}
{Targeted document(s) in .planning/codebase/}
{/gsd:scan -{}-focus quality}

\flagtable{%
\gsdflag{focus <area>} & Focus area: \texttt{tech}, \texttt{arch}, \texttt{quality}, \texttt{concerns}, \texttt{tech+arch} (default) \\
}

\commandentry{intel}
{Query, inspect, or refresh queryable codebase intelligence files. Requires \texttt{intel.enabled: true} in config.json.}
{/gsd:intel [query <term> | status | diff | refresh]}
{.planning/intel/ JSON files (stack, api-map, dependency-graph, file-roles, arch-decisions)}
{/gsd:intel query authentication}

\commandentry{onboard}
{Brownfield onboarding workflow that pairs map-codebase output with a new GSD project scaffold.}
{/gsd:onboard}
{PROJECT.md, REQUIREMENTS.md, ROADMAP.md tailored to the mapped codebase}
{/gsd:onboard}

\commandentry{migrate}
{Staged migration planner for legacy code — produces migration phases in the roadmap.}
{/gsd:migrate}
{Migration phases in ROADMAP.md}
{/gsd:migrate}

% ============================================================================
\chapter{AI Integration}

\commandentry{ai-integration-phase}
{AI framework selection wizard for integrating AI/LLM capabilities into a phase. Spawns 3 parallel specialist agents (domain researcher, framework selector, ai researcher, eval planner).}
{/gsd:ai-integration-phase [N]}
{\{phase\}-AI-SPEC.md with framework recommendation, implementation guidance, evaluation strategy}
{/gsd:ai-integration-phase 3}

\commandentry{eval-review}
{Retroactive audit of an implemented AI phase's evaluation coverage. Scores each eval dimension as COVERED/PARTIAL/MISSING.}
{/gsd:eval-review [N]}
{\{phase\}-EVAL-REVIEW.md with findings, gaps, and remediation guidance}
{/gsd:eval-review 3}

% ============================================================================
\chapter{Update Commands}

\commandentry{update}
{Update GSD with a changelog preview before applying.}
{/gsd:update}
{}
{/gsd:update}

\commandentry{reapply-patches}
{Restore local modifications after a GSD update. Merges back any local changes that the updater replaced.}
{/gsd:reapply-patches}
{}
{/gsd:reapply-patches}

% ============================================================================
\chapter{Code Quality}

\commandentry{code-review}
{Review source files changed during a phase for bugs, security vulnerabilities, and code quality problems. Three depth levels.}
{/gsd:code-review <N> [-{}-depth=quick|standard|deep] [-{}-files f1,f2,...]}
{\{phase\}-REVIEW.md with severity-classified findings}
{/gsd:code-review 2 -{}-depth=deep}

\flagtable{%
\gsdflag{depth=quick} & Pattern-matching only (\textasciitilde 2 min) \\
\gsdflag{depth=standard} & Per-file analysis with language-specific checks (default, 5--15 min) \\
\gsdflag{depth=deep} & Cross-file analysis including import graphs and call chains (15--30 min) \\
\gsdflag{files f1,f2} & Explicit file list; skips SUMMARY and git scoping \\
}

\commandentry{code-review-fix}
{Auto-fix issues found by code-review. Reads REVIEW.md, spawns a fixer agent, commits each fix atomically.}
{/gsd:code-review-fix <N> [-{}-all] [-{}-auto]}
{\{phase\}-REVIEW-FIX.md}
{/gsd:code-review-fix 3 -{}-auto}

\flagtable{%
\gsdflag{all} & Include Info findings (default: Critical + Warning only) \\
\gsdflag{auto} & Fix and re-review iteration loop, capped at 3 iterations \\
}

\commandentry{audit-fix}
{Autonomous audit-to-fix pipeline. Runs an audit, classifies findings, fixes auto-fixable issues with test verification, and commits each fix atomically.}
{/gsd:audit-fix [-{}-source <audit>] [-{}-severity high|medium|all] [-{}-max N] [-{}-dry-run]}
{Fix commits with test verification; classification report}
{/gsd:audit-fix -{}-severity high}

\flagtable{%
\gsdflag{source <audit>} & Which audit to run (default: \texttt{audit-uat}) \\
\gsdflag{severity} & Minimum severity to process: \texttt{high}, \texttt{medium}, \texttt{all} (default: \texttt{medium}) \\
\gsdflag{max N} & Maximum findings to fix (default: 5) \\
\gsdflag{dry-run} & Classify findings without fixing \\
}

\commandentry{secure-phase}
{Retroactively verify threat mitigations for a completed phase. Three operating modes: audit existing SECURITY.md, generate from PLAN.md threat model, or exit with guidance if phase not executed.}
{/gsd:secure-phase [N]}
{\{phase\}-SECURITY.md with threat verification results}
{/gsd:secure-phase 5}

\commandentry{docs-update}
{Generate or update project documentation verified against the codebase. Each doc writer explores the codebase directly — no hallucinated paths.}
{/gsd:docs-update [-{}-force] [-{}-verify-only]}
{Up to 9 doc files: README, architecture, API, getting-started, development, testing, configuration, deployment, contributing}
{/gsd:docs-update -{}-verify-only}

\flagtable{%
\gsdflag{force} & Skip preservation prompts, regenerate all docs \\
\gsdflag{verify-only} & Check existing docs for accuracy, no generation \\
}

% ============================================================================
\chapter{Fast \& Inline}

\commandentry{fast}
{Execute a trivial task inline — no subagents, no planning overhead. For typo fixes, config changes, small refactors, forgotten commits.}
{/gsd:fast [task description]}
{Direct edit and commit}
{/gsd:fast "fix typo in README"}

\begin{gsdnote}
\texttt{fast} is not a replacement for \gsdcolon{quick}. Use \gsdcolon{quick} for anything needing research, multi-step planning, or verification.
\end{gsdnote}

\commandentry{review}
{Cross-AI peer review of phase plans from external AI CLIs. Produces REVIEWS.md consumable by \gsdcolon{plan-phase -{}-reviews}.}
{/gsd:review -{}-phase N [-{}-gemini] [-{}-claude] [-{}-codex] [-{}-coderabbit] [-{}-opencode] [-{}-qwen] [-{}-cursor] [-{}-all]}
{\{phase\}-REVIEWS.md}
{/gsd:review -{}-phase 3 -{}-all}

\flagtable{%
\gsdflag{phase N} & Phase number to review (required) \\
\gsdflag{gemini} & Include Gemini CLI review \\
\gsdflag{claude} & Include Claude CLI review (separate session) \\
\gsdflag{codex} & Include Codex CLI review \\
\gsdflag{coderabbit} & Include CodeRabbit review \\
\gsdflag{opencode} & Include OpenCode review (via GitHub Copilot) \\
\gsdflag{qwen} & Include Qwen Code review \\
\gsdflag{cursor} & Include Cursor agent review \\
\gsdflag{all} & Include all available CLIs \\
}

\commandentry{pr-branch}
{Create a clean PR branch by filtering out \gsdfile{.planning/} commits so reviewers see only code changes.}
{/gsd:pr-branch [target-branch]}
{Clean git branch}
{/gsd:pr-branch develop}

% ============================================================================
\chapter{Backlog, Seeds \& Threads}

\commandentry{add-backlog}
{Add an idea to the backlog parking lot using 999.x numbering. Creates the phase directory immediately so discuss-phase and plan-phase work on it.}
{/gsd:add-backlog <description>}
{999.x phase directory}
{/gsd:add-backlog "GraphQL API layer"}

\commandentry{review-backlog}
{Review and promote backlog items to the active milestone. Per-item actions: Promote, Keep, Remove.}
{/gsd:review-backlog}
{Updated ROADMAP.md if any backlog promoted}
{/gsd:review-backlog}

\commandentry{plant-seed}
{Capture a forward-looking idea with trigger conditions — surfaces automatically at the right milestone. Solves context rot for ideas that don't belong to the current phase.}
{/gsd:plant-seed [idea summary]}
{.planning/seeds/SEED-NNN-slug.md}
{/gsd:plant-seed "Real-time collab when WebSocket infra is in"}

\commandentry{thread}
{Manage persistent context threads for cross-session work. Lighter weight than pause-work.}
{/gsd:thread [name | description]}
{.planning/threads/\{name\}.md}
{/gsd:thread fix-deploy-key-auth}

% ============================================================================
\chapter{State Management (CLI)}

State-management commands are invoked through the \gsdfile{gsd-tools.cjs} CLI rather than the workflow surface. They are still part of the v1.35.0 command inventory.

\commandentry{tools state validate}
{Detect drift between STATE.md and the actual filesystem.}
{node gsd-tools.cjs state validate}
{Validation report showing any drift}
{node gsd-tools.cjs state validate}

\commandentry{tools state sync}
{Reconstruct STATE.md from actual project state on disk.}
{node gsd-tools.cjs state sync [-{}-verify]}
{Updated STATE.md}
{node gsd-tools.cjs state sync -{}-verify}

\flagtable{%
\gsdflag{verify} & Dry-run mode — show proposed changes without writing \\
}

\commandentry{tools state planned-phase}
{Record state transition after plan-phase completes (Planned / Ready to execute).}
{node gsd-tools.cjs state planned-phase -{}-phase N -{}-plans N}
{Updated STATE.md with post-planning state}
{node gsd-tools.cjs state planned-phase -{}-phase 3 -{}-plans 2}

% ============================================================================
\chapter{Community}

\commandentry{join-discord}
{Open the Discord community invite.}
{/gsd:join-discord}
{}
{/gsd:join-discord}

\section*{Community Hooks}

Optional git and session hooks gated behind \texttt{hooks.community: true} in \gsdfile{.planning/config.json}. All are no-ops unless explicitly enabled.

\begin{tabularx}{\linewidth}{@{}l X@{}}
\toprule
\textbf{\footnotesize Hook} & \textbf{\footnotesize Purpose} \\
\midrule
\gsdfile{gsd-validate-commit.sh} & Enforce Conventional Commits format on git commit messages \\
\gsdfile{gsd-session-state.sh} & Track session state transitions \\
\gsdfile{gsd-phase-boundary.sh} & Enforce phase boundary checks \\
\bottomrule
\end{tabularx}

\smallskip
\noindent Enable with:
\begin{lstlisting}
{ "hooks": { "community": true } }
\end{lstlisting}

% ============================================================================
\chapter{Configuration Reference}

GSD stores project settings in \gsdfile{.planning/config.json}. The file is created by \gsdcolon{new-project} and can be updated interactively via \gsdcolon{settings} or per-key via \texttt{node gsd-tools.cjs config-set <key> <value>}.

\section{Core Settings}

\begin{longtable}{@{}l l l@{}}
\toprule
\textbf{\footnotesize Setting} & \textbf{\footnotesize Type} & \textbf{\footnotesize Default} \\
\midrule
\endhead
\texttt{mode} & \texttt{interactive} \textbar{} \texttt{yolo} & \texttt{interactive} \\
\texttt{granularity} & \texttt{coarse} \textbar{} \texttt{standard} \textbar{} \texttt{fine} & \texttt{standard} \\
\texttt{model\_profile} & \texttt{quality/balanced/budget/inherit} & \texttt{balanced} \\
\texttt{project\_code} & string prefix for phase dirs & (none) \\
\texttt{response\_language} & language code (\texttt{pt}, \texttt{ko}, \texttt{ja}) & (none) \\
\texttt{context\_profile} & \texttt{dev} \textbar{} \texttt{research} \textbar{} \texttt{review} & (none) \\
\texttt{claude\_md\_path} & custom CLAUDE.md output path & (none) \\
\texttt{security\_enforcement} & boolean & \texttt{true} \\
\texttt{security\_asvs\_level} & 1 \textbar{} 2 \textbar{} 3 & \texttt{1} \\
\texttt{security\_block\_on} & \texttt{high} \textbar{} \texttt{medium} \textbar{} \texttt{low} & \texttt{high} \\
\bottomrule
\end{longtable}

\section{Workflow Toggles}

All workflow toggles follow the \textbf{absent = enabled} pattern. Missing keys default to \texttt{true} unless otherwise noted.

\begin{longtable}{@{}l l l@{}}
\toprule
\textbf{\footnotesize Setting} & \textbf{\footnotesize Type} & \textbf{\footnotesize Default} \\
\midrule
\endhead
\texttt{workflow.research} & boolean & \texttt{true} \\
\texttt{workflow.plan\_check} & boolean & \texttt{true} \\
\texttt{workflow.verifier} & boolean & \texttt{true} \\
\texttt{workflow.auto\_advance} & boolean & \texttt{false} \\
\texttt{workflow.nyquist\_validation} & boolean & \texttt{true} \\
\texttt{workflow.ui\_phase} & boolean & \texttt{true} \\
\texttt{workflow.ui\_safety\_gate} & boolean & \texttt{true} \\
\texttt{workflow.node\_repair} & boolean & \texttt{true} \\
\texttt{workflow.node\_repair\_budget} & number & \texttt{2} \\
\texttt{workflow.research\_before\_questions} & boolean & \texttt{false} \\
\texttt{workflow.discuss\_mode} & \texttt{discuss} \textbar{} \texttt{assumptions} & \texttt{discuss} \\
\texttt{workflow.skip\_discuss} & boolean & \texttt{false} \\
\texttt{workflow.text\_mode} & boolean & \texttt{false} \\
\texttt{workflow.use\_worktrees} & boolean & \texttt{true} \\
\texttt{workflow.code\_review} & boolean & \texttt{true} \\
\texttt{workflow.code\_review\_depth} & \texttt{quick/standard/deep} & \texttt{standard} \\
\texttt{workflow.plan\_bounce} & boolean & \texttt{false} \\
\texttt{workflow.plan\_bounce\_script} & string path & (none) \\
\texttt{workflow.plan\_bounce\_passes} & number & \texttt{2} \\
\texttt{workflow.code\_review\_command} & shell command & (none) \\
\texttt{workflow.tdd\_mode} & boolean & \texttt{false} \\
\texttt{workflow.cross\_ai\_execution} & boolean & \texttt{false} \\
\texttt{workflow.cross\_ai\_command} & shell command & (none) \\
\texttt{workflow.cross\_ai\_timeout} & seconds & \texttt{300} \\
\bottomrule
\end{longtable}

\section{Planning, Hooks \& Features}

\begin{longtable}{@{}l l l@{}}
\toprule
\textbf{\footnotesize Setting} & \textbf{\footnotesize Type} & \textbf{\footnotesize Default} \\
\midrule
\endhead
\texttt{planning.commit\_docs} & boolean & \texttt{true} \\
\texttt{planning.search\_gitignored} & boolean & \texttt{false} \\
\texttt{hooks.context\_warnings} & boolean & \texttt{true} \\
\texttt{hooks.workflow\_guard} & boolean & \texttt{false} \\
\texttt{hooks.community} & boolean & \texttt{false} \\
\texttt{features.thinking\_partner} & boolean & \texttt{false} \\
\texttt{features.global\_learnings} & boolean & \texttt{false} \\
\texttt{intel.enabled} & boolean & \texttt{false} \\
\texttt{learnings.max\_inject} & number & \texttt{10} \\
\texttt{agent\_skills} & object & \texttt{\{\}} \\
\bottomrule
\end{longtable}

\begin{gsdnote}
If \gsdfile{.planning/} is in \gsdfile{.gitignore}, \texttt{planning.commit\_docs} is automatically forced to \texttt{false} regardless of config.json. This prevents git errors.
\end{gsdnote}

\section{Parallelization}

\begin{longtable}{@{}l l l@{}}
\toprule
\textbf{\footnotesize Setting} & \textbf{\footnotesize Type} & \textbf{\footnotesize Default} \\
\midrule
\endhead
\texttt{parallelization.enabled} & boolean & \texttt{true} \\
\texttt{parallelization.plan\_level} & boolean & \texttt{true} \\
\texttt{parallelization.task\_level} & boolean & \texttt{false} \\
\texttt{parallelization.skip\_checkpoints} & boolean & \texttt{true} \\
\texttt{parallelization.max\_concurrent\_agents} & number & \texttt{3} \\
\texttt{parallelization.min\_plans\_for\_parallel} & number & \texttt{2} \\
\bottomrule
\end{longtable}

\begin{gsdwarning}
When parallelization is enabled, executor agents commit with \gsdflag{no-verify} to avoid build-lock contention (e.g., cargo lock fights in Rust projects). The orchestrator validates hooks once after each wave completes. Set \texttt{parallelization.enabled: false} if you need hooks to run per-commit.
\end{gsdwarning}

\section{Git Branching}

\begin{longtable}{@{}l l l@{}}
\toprule
\textbf{\footnotesize Setting} & \textbf{\footnotesize Type} & \textbf{\footnotesize Default} \\
\midrule
\endhead
\texttt{git.branching\_strategy} & \texttt{none} \textbar{} \texttt{phase} \textbar{} \texttt{milestone} & \texttt{none} \\
\texttt{git.phase\_branch\_template} & string template & \texttt{gsd/phase-\{phase\}-\{slug\}} \\
\texttt{git.milestone\_branch\_template} & string template & \texttt{gsd/\{milestone\}-\{slug\}} \\
\texttt{git.quick\_branch\_template} & string template \textbar{} null & \texttt{null} \\
\bottomrule
\end{longtable}

Template variables: \texttt{\{phase\}} (zero-padded), \texttt{\{slug\}} (lowercase, hyphenated), \texttt{\{milestone\}} (e.g., \texttt{v1.0}), \texttt{\{num\}} / \texttt{\{quick\}} (quick task ID).

\section{Gates, Safety \& Review}

\begin{longtable}{@{}l l l@{}}
\toprule
\textbf{\footnotesize Setting} & \textbf{\footnotesize Type} & \textbf{\footnotesize Default} \\
\midrule
\endhead
\texttt{gates.confirm\_project} & boolean & \texttt{true} \\
\texttt{gates.confirm\_phases} & boolean & \texttt{true} \\
\texttt{gates.confirm\_roadmap} & boolean & \texttt{true} \\
\texttt{gates.confirm\_breakdown} & boolean & \texttt{true} \\
\texttt{gates.confirm\_plan} & boolean & \texttt{true} \\
\texttt{gates.execute\_next\_plan} & boolean & \texttt{true} \\
\texttt{gates.issues\_review} & boolean & \texttt{true} \\
\texttt{gates.confirm\_transition} & boolean & \texttt{true} \\
\texttt{safety.always\_confirm\_destructive} & boolean & \texttt{true} \\
\texttt{safety.always\_confirm\_external\_services} & boolean & \texttt{true} \\
\texttt{review.models.gemini} & string & (CLI default) \\
\texttt{review.models.claude} & string & (CLI default) \\
\texttt{review.models.codex} & string & (CLI default) \\
\texttt{review.models.opencode} & string & (CLI default) \\
\texttt{review.models.qwen} & string & (CLI default) \\
\texttt{review.models.cursor} & string & (CLI default) \\
\bottomrule
\end{longtable}

\section{Manager Passthrough Flags}

\begin{longtable}{@{}l l l@{}}
\toprule
\textbf{\footnotesize Setting} & \textbf{\footnotesize Type} & \textbf{\footnotesize Default} \\
\midrule
\endhead
\texttt{manager.flags.discuss} & string & (none) \\
\texttt{manager.flags.plan} & string & (none) \\
\texttt{manager.flags.execute} & string & (none) \\
\bottomrule
\end{longtable}

Flags are appended to each dispatched command. Invalid flag tokens are sanitized and logged as warnings; only recognized GSD flags are passed through.

\section{Model Profiles}

\begin{longtable}{@{}l l l l l@{}}
\toprule
\textbf{\footnotesize Agent} & \textbf{\footnotesize quality} & \textbf{\footnotesize balanced} & \textbf{\footnotesize budget} & \textbf{\footnotesize inherit} \\
\midrule
\endhead
gsd-planner & Opus & Opus & Sonnet & inherit \\
gsd-roadmapper & Opus & Sonnet & Sonnet & inherit \\
gsd-executor & Opus & Sonnet & Sonnet & inherit \\
gsd-phase-researcher & Opus & Sonnet & Haiku & inherit \\
gsd-project-researcher & Opus & Sonnet & Haiku & inherit \\
gsd-research-synthesizer & Sonnet & Sonnet & Haiku & inherit \\
gsd-debugger & Opus & Sonnet & Sonnet & inherit \\
gsd-codebase-mapper & Sonnet & Haiku & Haiku & inherit \\
gsd-verifier & Sonnet & Sonnet & Haiku & inherit \\
gsd-plan-checker & Sonnet & Sonnet & Haiku & inherit \\
gsd-integration-checker & Sonnet & Sonnet & Haiku & inherit \\
gsd-nyquist-auditor & Sonnet & Sonnet & Haiku & inherit \\
\bottomrule
\end{longtable}

Override specific agents via \texttt{model\_overrides}. Valid values: \texttt{opus}, \texttt{sonnet}, \texttt{haiku}, \texttt{inherit}, or any fully-qualified model ID (\texttt{openai/o3}, \texttt{google/gemini-2.5-pro}, etc.). Non-Claude runtimes ship with \texttt{resolve\_model\_ids: "omit"} so each agent inherits the runtime's default model unless overridden.

\section{Environment Variables}

\begin{tabularx}{\linewidth}{@{}l X@{}}
\toprule
\textbf{\footnotesize Variable} & \textbf{\footnotesize Purpose} \\
\midrule
\texttt{CLAUDE\_CONFIG\_DIR} & Override default config directory (\gsdfile{\textasciitilde/.claude/}) \\
\texttt{GEMINI\_API\_KEY} & Detected by context monitor to switch hook event name \\
\texttt{WSL\_DISTRO\_NAME} & Detected by installer for WSL path handling \\
\texttt{GSD\_SKIP\_SCHEMA\_CHECK} & Skip schema drift detection during execute-phase \\
\texttt{GSD\_PROJECT} & Override project root for multi-project workspace support \\
\bottomrule
\end{tabularx}

\section{Global Defaults}

Save settings as global defaults for future projects at \gsdfile{\textasciitilde/.gsd/defaults.json}. When \gsdcolon{new-project} creates a new config.json, it reads global defaults and merges them as the starting configuration. Per-project settings always override globals.

% ============================================================================
\chapter{File Manifest}

GSD artifacts live under \gsdfile{.planning/} at the project root (or at \gsdfile{\$GSD\_PROJECT/.planning/} if overridden). The canonical layout for a typical v1.35.0 project:

\begin{lstlisting}
.planning/
  PROJECT.md                  # Project vision, always loaded
  REQUIREMENTS.md             # Scoped v1/v2 requirements with phase traceability
  ROADMAP.md                  # Phase list, status, dependencies
  STATE.md                    # Decisions, blockers, position across sessions
  config.json                 # Workflow, model, git, gate, safety configuration
  MILESTONES.md               # Archive of completed milestones
  CLAUDE.md                   # Generated; written at project root by default
  continue-here.md            # Written by /gsd:pause-work
  research/                   # Project-level ecosystem research
  codebase/                   # /gsd:map-codebase and /gsd:scan output
  intel/                      # /gsd:intel JSON files (stack, api-map, etc.)
  seeds/
    SEED-NNN-slug.md          # /gsd:plant-seed artifacts
  threads/
    {thread-name}.md          # /gsd:thread cross-session knowledge stores
  todos/                      # Captured ideas from /gsd:add-todo and /gsd:note
  quick/
    NNN-slug/                 # /gsd:quick ad-hoc task directory
      PLAN.md
      SUMMARY.md
  forensics/
    report-{timestamp}.md     # /gsd:forensics output
  learnings/
    {phase}-LEARNINGS.md      # /gsd:extract-learnings output
  reports/
    SESSION_REPORT.md         # /gsd:session-report
    MILESTONE_SUMMARY-v{v}.md # /gsd:milestone-summary
  ui-reviews/                 # Screenshots from /gsd:ui-review
  phases/
    NN-slug/                  # One directory per phase
      CONTEXT.md              # From /gsd:discuss-phase
      DISCUSSION-LOG.md       # Audit trail for discuss-phase
      RESEARCH.md             # From /gsd:plan-phase (or /gsd:research-phase)
      PLAN.md                 # Aggregate phase plan
      NN-PLAN.md              # Per-wave plans (NN = wave number)
      VALIDATION.md           # Nyquist validation data
      NN-SUMMARY.md           # Per-plan execution summary
      VERIFICATION.md         # From /gsd:verify-work and executor wrap-up
      UAT.md                  # User acceptance test results
      UI-SPEC.md              # From /gsd:ui-phase
      UI-REVIEW.md            # From /gsd:ui-review
      REVIEW.md               # From /gsd:code-review
      REVIEW-FIX.md           # From /gsd:code-review-fix
      REVIEWS.md              # From /gsd:review (cross-AI)
      SECURITY.md             # From /gsd:secure-phase
      AI-SPEC.md              # From /gsd:ai-integration-phase
      EVAL-REVIEW.md          # From /gsd:eval-review
\end{lstlisting}

\section*{Workspace Layout}

Workspaces created by \gsdcolon{new-workspace} live under \gsdfile{\textasciitilde/gsd-workspaces/<name>/} by default, with their own \gsdfile{.planning/} directory, a \gsdfile{WORKSPACE.md} manifest, and one or more repo copies (as worktrees or clones).

\section*{Global Layout}

\begin{lstlisting}
~/.gsd/
  defaults.json               # Global defaults merged into new projects
~/.claude/                    # Runtime config dir (overridable via CLAUDE_CONFIG_DIR)
\end{lstlisting}

\vspace{1em}
\begin{center}
\small\textit{End of reference. GSD v1.35.0. This document is CC BY-SA 4.0 — The Tibsfox Team.}
\end{center}

\end{document}
